% Part:sets-functions-relations
% Chapter: size-of-sets
% Section: introduction

\documentclass[../../../include/open-logic-section]{subfiles}

\begin{document}

\olfileid{sfr}{siz}{int}

\olsection{ভূমিকা}

জর্জ কান্টর ১৮৭০-এর দশকে যখন সেটতত্ত্ব গড়ে তোলেন, তখন তাঁর অন্যতম উদ্দেশ্য ছিল একটি অসীম সমষ্টির ধারণাকে গ্রহণযোগ্য করে তোলা—মধ্যযুগীয় চিন্তকেরা যাকে বাস্তবায়িত অসীম বলতেন। বিভিন্ন সেটের \emph{আকার} নিয়ে তাঁর আলোচনা ছিল এর একটি প্রধান অংশ। যদি $a$, $b$ ও $c$ সকলেই ভিন্ন হয়, তবে স্বাভাবিক বোধে সেট $\{a, b, c\}$, $\{a, b\}$-এর চেয়ে \emph{বড়}। কিন্তু অসীম সেটগুলির ক্ষেত্রে কী হবে? সেগুলি কি সকলেই পরস্পরের সমান বড়? দেখা যায়, তা নয়।

এখানে প্রথম গুরুত্বপূর্ণ ধারণা হলো তালিকায়ন। একটি সসীম সেটের সকল !!{element}s লিখে আমরা সেটিকে তালিকায় প্রকাশ করতে পারি। তালিকাটিকেই অসীম হতে দিলে কিছু অসীম সেটের সকল !!{element}s-ও আমরা তালিকায় লিখতে পারি। এই ধরনের সেটকে !!{enumerable} বলা হয়। কান্টরের বিস্ময়কর ফল ছিল এই যে কিছু অসীম সেট !!{enumerable} নয়; এই অধ্যায়ের শেষে আমরা ফলটি সম্পূর্ণভাবে বুঝতে পারব।

\end{document}
